Mean (arithmetic average):  \bar{x} = \frac{x_1 + x_2 + \cdot\cdot\cdot + x_n}{n} = \frac{\sum x_i}{n}

Weighted mean:  \bar{x} = \frac{n_1\bar{x}_1 + n_2\bar{x}_2}{n_1 + n_2}

Example (2020 \#5): 
\bar{x}_1 = 32 over 50 numbers,  \bar{x}_2 = 53  over 70 numbers
\bar{x} = \frac{50(32) + 70(53)}{120} = \frac{1600 + 3710}{120} = \frac{5310}{120} = \frac{177}{4} = 44.25

Median: middle value when data is ordered.
For  n  values: position = \frac{n+1}{2}
If  n  even: average of the two middle values.

Example scores:  62,68,73,77,78,82,86,86,90,92  (n=10)
Median = \frac{78+82}{2} = 80

Mode: most frequently occurring value(s).

Median adjustment (2026 \#19):
Each score increased by  \frac{1}{10} \times 80 = 8
Fran's score (lowest: 62):  62 + 8 = 70
Percentage increase:  \frac{8}{62} \times 100 \approx 12.903\%

Geometric mean of  a  and  b:  G = \sqrt{ab}

Harmonic mean of  a  and  b:  H = \frac{2ab}{a+b}

Example (2022 \#16): harmonic mean of 10 and 24:
H = \frac{2(10)(24)}{10+24} = \frac{480}{34} = \frac{240}{17}

Average speed: use  \frac{total distance}{total time}

Example (2024 \#13):  \frac{1}{3}d  at 70 km/h,  \frac{2}{3}d  at 35 km/h
t = \frac{d/3}{70} + \frac{2d/3}{35} = \frac{d}{210} + \frac{2d}{105} = \frac{5d}{210}
\bar{v} = \frac{d}{5d/210} = \frac{210}{5} = 42  km/h

Mean of  S = \left\{\frac{1}{2}, \frac{1}{3}, \frac{1}{4}, \frac{1}{5}\right\}:
\bar{x} = \frac{1}{4}\!\left(\frac{1}{2}+\frac{1}{3}+\frac{1}{4}+\frac{1}{5}\right) = \frac{1}{4} \cdot \frac{30+20+15+12}{60} = \frac{77}{240}
Median:  \frac{1/4 + 1/3}{2} = \frac{7}{48}
